\chapter[chap:phase5]{Phase 5: Netzwerk und Hub}

\section{Ziel}

Repositories werden zwischen Klonen und einem Hub synchronisiert.
Der Hub setzt die Berechtigungen durch, zusätzlich zur lokalen Durchsetzung aus
\in{Kapitel}[chap:phase4].
Was er dabei hinzufügt, ist eine Ordnung, die jemand schreibt, der nicht selbst der Akteur ist.

\section{Umfang}

\startitemize
\item Remotes, \type{clone}, \type{fetch} und \type{push} im Kern, hinter einer Schnittstelle für
den Transport (ADR-022).
\item Ein Übertragungsprotokoll über HTTP, das ganze Objekte als Strom überträgt (ADR-023).
\item Anmeldung am Hub mit dem Schlüssel aus dem Akteursregister (ADR-024).
\item Remote-Referenzen und das signierte Ref-Update mit seiner Aufzeichnung auf dem Hub
(ADR-008, ADR-025).
\item Durchsetzung von Schreibhoheit, Merge-Gate und Eskalation auf dem Hub, mit denselben
Operationen des Kerns wie lokal (ADR-025).
\item Der Anker des Operation Log auf dem Hub und der Schlüssel des Hubs (ADR-026).
\item Namespaces für Organisationen, Kurse und Projekte, mit einer Namespace-Policy (ADR-027).
\item Der Schlüssel auf einem Hardware-Token als zweiter Schlüsseltyp des Registers (ADR-020).
\stopitemize

Die Phase hat begonnen: Die Fragen des Entwurfs sind entschieden und in den ADRs festgehalten,
umgesetzt ist noch nichts.
Dieses Kapitel beschreibt den entschiedenen Entwurf; jeder Abschnitt wächst mit dem Teil, der
gebaut wird, und die Befehle kommen mit ihrer Umsetzung in \in{Kapitel}[chap:cli].
Der Hub selbst gehört zum getrennten Projekt stemma-hub (ADR-003).
Das Kapitel beschreibt deshalb das Protokoll und die Regeln, die ein Klon kennt und nachrechnet;
wie der Hub seine Daten ablegt und betrieben wird, steht hier nicht.

\section[sec:remote-split]{Klon, Remote und Hub}

Phase 4 setzt die Regeln auf der Maschine des Akteurs durch, für die Prozesse, die durch den Kern
gehen.
Ein Hub ist ein zweiter Ort, der dieselben Beweise noch einmal prüft, bevor sich eine geteilte
Referenz bewegt.
Damit beide Orte zum selben Urteil kommen, gibt es die Prüfungen nur einmal, im Kern.

Der Schnitt auf der Seite des Klons verläuft entlang der Ein- und Ausgabe.
Zustand und Logik liegen im Kern: die Konfiguration der Remotes, die Remote-Referenzen, die
Berechnung der fehlenden Objekte, die Prüfung empfangener Objekte und die Einträge von
\type{fetch} und \type{push} im Operation Log.
Der Kern spricht mit der Gegenseite über eine Schnittstelle (\type{Transport}), die nur Nachrichten
überträgt und nichts über HTTP annimmt.
Ein eigenes Crate (\type{stemma-remote}) enthält den HTTP-Client, der sie umsetzt; der Kern selbst
bekommt keine Netzwerk-Abhängigkeit.

Die Nachrichten, auf die sich Klon und Hub einigen müssen, sind Typen des Kerns
(\type{stemma_core::protocol}), mit einem Schreiber und einem strengen Leser für beide Seiten.
Jede Nachricht beginnt mit der Zeile \type{stemma-protocol <n>}.
Diese Protokollversion ist unabhängig von der Version des Repository-Formats (ADR-009) und von
der Version des Programms; eine Änderung braucht ein ADR und eine neue Nummer.
Ein Hub kann mehrere Versionen bedienen, ein Klon spricht eine.

Der Hub öffnet jedes Repository als Bare-Repository, also ohne Arbeitskopie, und ruft die
Operationen des Kerns auf.
Eine einzige Operation nimmt ein Ref-Update an und legt die Reihenfolge der Prüfungen fest
(\in{Abschnitt}[sec:hub-enforcement]).
Der Hub liefert ihr die Belege und wendet ihr Ergebnis an; er setzt keine Prüfung selbst
zusammen und kann deshalb keine auslassen.

\startinfobox
{\bf Wie Klon und Hub gemeinsam getestet werden.}
Die Tests dieses Repositorys verbinden zwei Repositories im Speicher über \type{Transport}; die
Gegenseite ist dabei die annehmende Operation des Kerns, kein zweiter Server.
Dazu kommen eingecheckte Beispielnachrichten, gegen die auch der Hub seine Leser prüft.
Tests mit echtem HTTP liegen beim Hub.
\stopinfobox

\section[sec:transfer]{Das Übertragungsprotokoll}

Ein Klon und ein Hub müssen herausfinden, welche Objekte der anderen Seite fehlen, und sie
übertragen.
Der Entwurf nutzt dafür eine Voraussetzung aus ADR-015: Jeder Klon hält den vollen
Objektspeicher.
Wer einen Commit hat, hat alles, was von ihm aus erreichbar ist.

\subsection{Die Einheit der Übertragung}

Eine Anfrage überträgt einen Objektstrom: ganze Objekte in der Kodierung des Speichers
(\type{<Art> <Länge>}, ein NUL-Byte, der Inhalt; \in{Abschnitt}[sec:object-ids]), ohne Deltas.
Der Empfänger legt sie über den bestehenden Schreibweg ab, es entsteht also kein zweites
Speicherformat.
Das Objekt bleibt die Einheit, die ihr Hash adressiert, aber es ist nicht die Einheit einer
Anfrage.
Ein Objekt je Anfrage würde für einen Klon dieses Repositorys etwa 4500 Anfragen kosten, bei
einem Verlauf von einigen Gigabyte $10^5$ bis $10^6$.

Der ganze Strom wird komprimiert, als Eigenschaft des Transports und mit zstd (ADR-029); der
Speicher bleibt, wie ADR-007 ihn festlegt.
Eine Messung am Git-Verlauf dieses Repositorys gab die Richtung: 4483 Objekte sind roh
103,3 MB, mit zstd und einem Fenster über den ganzen Verlauf 2,4 MB, als Git-Pack mit Deltas
7,4 MB.
Das lange Fenster findet die Wiederholungen zwischen den Versionen einer Datei selbst; deshalb
kommt das Format ohne Deltas aus.
Die Zahlen stammen von einem einzigen Verlauf und sind kein Schwellenwert.

\subsection{Fetch und Klon}

\startitemize[n]
\item Der Klon liest die Liste aller Referenzen des Hubs in einer Anfrage.
\item Er sendet in einer zweiten Anfrage die Commits, die er will (\type{want}), und die Commits,
die er hat (\type{have}): die Spitzen seiner Branches und seiner Remote-Referenzen.
\item Der Hub läuft den Graphen ab und antwortet mit dem Strom aller Objekte, die von den
\type{want}-Commits aus erreichbar sind und von den bekannten \type{have}-Commits aus nicht.
\item Der Klon prüft, was er empfangen hat, und bewegt erst danach seine Remote-Referenzen.
\stopitemize

Ein Klon ist dieselbe Anfrage ohne \type{have}.
Eine zweite Runde gibt es nicht.
Ein \type{have}, das der Hub nicht kennt, übergeht er; er sendet dann zu viel, nie zu wenig, und
ein doppeltes Objekt schadet nicht, weil sein Hash es adressiert.
Fetch und Klon eines öffentlichen Repositorys kosten damit zwei Anfragen, unabhängig von der
Grösse des Verlaufs.

Ein \type{want} ist an eine Referenz gebunden: Der Hub liefert nur, was von einer Referenz eines
Repositorys aus erreichbar ist, das der Anfragende lesen darf.
Kein Endpunkt liefert ein Objekt allein über seine ID, auch nicht in einem öffentlichen
Repository.
Sonst bekäme jeder, der eine ID kennt, auch die Objekte eines abgelehnten Push.

\subsection{Push}

\startitemize[n]
\item Der Klon berechnet, was dem Hub fehlt: alles, was von der neuen Spitze aus erreichbar ist
und von seinen Remote-Referenzen aus nicht.
Er fragt den Hub nicht vorher.
\item Er lädt diese Objekte als Strom hoch.
Der Hub berechnet ihre Hashes, legt sie unbenannt ab und bestätigt, sobald sie dauerhaft
gespeichert sind.
\item Er sendet ein kleines, signiertes Ref-Update (\in{Abschnitt}[sec:ref-update]).
Erst jetzt prüft der Hub, ob die Objekte zusammenhängen, ruft die annehmende Operation auf und
bewegt die Referenz mit Compare-and-Swap.
\stopitemize

Der Upload ist idempotent: Er kann wiederholt und an Commit-Grenzen in mehrere Anfragen geteilt
werden.
Ein abgelehnter Push hinterlässt unbenannte Objekte auf dem Hub, die er nie ausliefert und später
entfernt.
Ist die Remote-Referenz des Klons veraltet, scheitert das Compare-and-Swap (\type{Stale}), und der
Klon holt zuerst den neuen Stand.

\subsection{Was der Empfänger prüft}

Beide Seiten prüfen, der Klon beim Fetch und der Hub beim Push.
Der Empfänger wiederholt den Lauf des Senders: von den neuen Spitzen abwärts bis zu Commits, die
von seinen eigenen Referenzen aus erreichbar sind.
Jedes Objekt auf diesem Weg muss vorhanden sein, die richtige Art haben und vom strengen Leser
gelesen werden (\in{Abschnitt}[sec:strict-readers]).
Der Rahmen eines Objekts im Strom trägt keine ID.
Der Empfänger berechnet sie selbst; ein beschädigtes Objekt ist deshalb ein anderes Objekt, und
der Lauf meldet das erwartete als fehlend.
Ein Objekt, das der Lauf nicht erreicht, bleibt unbenannt im Speicher und ist kein Fehler.

Die Invariante dahinter: Eine Referenz zeigt nur auf einen Commit, dessen ganzer Verlauf im
Speicher liegt und lesbar ist.
Sie hält, weil die Referenz sich erst nach dem Lauf bewegt, und sie ist die Voraussetzung dafür,
dass eine einzige Runde mit \type{want} und \type{have} genügt.

\subsection{Unterbrochene Übertragung}

Die Reihenfolge des Stroms gehört zum Format: Vorfahren zuerst, und für jeden Commit seine neuen
Blobs und Trees vor dem Commit selbst.
Jeder Anfang des Stroms, der mit einem Commit endet, ist damit ein vollständiger Stand.
Der Klon prüft Commit für Commit und hält in \type{.stemma} fest, welche Commits vollständig
angekommen sind.
Nach einem Abbruch sendet er dieselbe Anfrage mit diesen Commits als weitere \type{have}-Zeilen.
Die Wiederaufnahme ist ein gewöhnlicher Fetch, und der Hub hält keinen Zustand.
Eine Grenze bleibt: Ein einzelner Commit, der grösser ist, als eine Verbindung trägt, lässt sich
nicht wiederaufnehmen.

\subsection{Nachrichten und Anfragen}

Die kleinen Nachrichten sind zeilenweiser Text in fester Reihenfolge, wie die Einträge des
Operation Log; der Objektstrom ist binär und beginnt mit derselben ersten Zeile.
Es gibt keine Listen von Fähigkeiten und keine optionalen Felder: Jede Erweiterung ist eine neue
Protokollversion.
Die Adresse eines Remote ist die Adresse des Repositorys,
\type{<Hub>/<Namespace>/<Repository>}, und alle Endpunkte liegen darunter
(\in{Tabelle}[tab:hub-requests]).

\startplacetable[title={Die Anfragen eines Klons an ein Repository auf dem Hub},
reference=tab:hub-requests]
\startxtable[frame=off,option=stretch]
\startxrow[topframe=on,bottomframe=on]
\startxcell[width=0.3\textwidth] {\bf Anfrage} \stopxcell
\startxcell {\bf Bedeutung} \stopxcell
\stopxrow
\startxrow
\startxcell \type{POST challenge} \stopxcell
\startxcell Der Hub antwortet mit einer Challenge. \stopxcell
\stopxrow
\startxrow
\startxcell \type{POST token} \stopxcell
\startxcell Die signierte Antwort des Klons; der Hub antwortet mit einem Token. \stopxcell
\stopxrow
\startxrow
\startxcell \type{GET refs} \stopxcell
\startxcell Die Liste der Referenzen, mit der ID des Repositorys und dem Schlüssel des Hubs.
\stopxcell
\stopxrow
\startxrow
\startxcell \type{POST fetch} \stopxcell
\startxcell \type{want} und \type{have}; die Antwort ist der Objektstrom. \stopxcell
\stopxrow
\startxrow
\startxcell \type{POST objects} \stopxcell
\startxcell Der Upload eines Objektstroms. \stopxcell
\stopxrow
\startxrow
\startxcell \type{POST refs} \stopxcell
\startxcell Das Ref-Update; die Antwort ist die Quittung. \stopxcell
\stopxrow
\startxrow[bottomframe=on]
\startxcell \type{POST anchors} \stopxcell
\startxcell Ein Anker des Operation Log; die Antwort ist die Bestätigung. \stopxcell
\stopxrow
\stopxtable
\stopplacetable

Weitere Anfragen lesen die Aufzeichnung, die signierten Erklärungen und die Anker und reichen
Erklärungen ein; ADR-023 führt sie auf.
Der Statuscode nennt nur die grobe Klasse eines Fehlers, der genaue Grund steht als Nachricht im
Rumpf der Antwort.
Der Klon folgt keiner Weiterleitung, damit ein Token nie an einen anderen Host geht.
Ein Fetch oder Klon eines öffentlichen Repositorys kostet 2 Anfragen, der eines privaten 4, ein
Push 5: die Anmeldung mit zwei Anfragen, der Upload, das Ref-Update und der Anker.
Ein Fetch liest ausserdem die Erklärungen offener Pull Requests und die Namespace-Policy, was je
eine Anfrage hinzufügen kann.

\startinfobox
{\bf Wann ein Pack-Format nötig wird.}
ADR-023 nennt die Auslöser, ohne Zahlen: Die Kompression des Stroms trägt nicht mehr, weil der
Verlauf grösser ist als das Fenster; der Hub wird zum Engpass eines Klons; ein einzelner Commit ist
grösser, als eine Verbindung trägt; oder der Speicher selbst braucht Packs (ADR-007).
Die Schwellen liefert eine Messung an echten Stemma-Speichern.
In jedem Fall ändert sich keine Objekt-ID.
\stopinfobox

\section[sec:hub-auth]{Anmeldung am Hub}

Der Hub muss wissen, welcher Schlüssel eine Anfrage sendet.
Phase 4 kennt dafür nur den Beweis: eine Signatur mit dem Schlüssel, den das Register dem Akteur
zuordnet (\in{Abschnitt}[sec:signatures]).
Die Anmeldung am Hub verwendet denselben Schlüssel und führt kein zweites Geheimnis ein.

\startitemize[n]
\item Der Hub gibt eine Challenge aus, einen Wert (Nonce), den er nur für Sekunden annimmt.
\item Der Schlüssel signiert sie, einmal je Befehl.
Die Signatur ist ein eigener Typ von Beweis und bindet die Adresse des Hubs, wie der Klon sie
aufruft, die Challenge und den Schlüssel.
\item Der Hub antwortet mit einem Token für die folgenden Anfragen.
\stopitemize

Das Token ist kurzlebig und kein Beweis.
Es gelangt in kein Repository, wird nicht auf die Platte geschrieben und sagt nur, welcher
Schlüssel spricht.
Was dieser Schlüssel darf, prüft der Hub bei jeder Anfrage gegen das Register an der Spitze von
\type{main}; ein Widerruf wirkt deshalb sofort.
Die Bindung an die Adresse verhindert, dass ein fremder Hub eine Challenge weiterreicht.
TLS ist für jede Adresse Pflicht, die ein Netz überquert (ADR-028).

Auf dem Hub hat ein Mensch ein Konto: einen Namen mit öffentlichen Schlüsseln, ohne eigenes
Geheimnis.
Das Konto und der Eintrag im Register sind über denselben Schlüssel verbunden, eine Tabelle, die
nur der Hub kennt, gibt es nicht.
Für alles, was ein Repository ändert, zählt allein das Register.
Das Konto braucht es nur oberhalb des Repositorys: um eines anzulegen, für einen Namespace und um
ohne Eintrag im Register zu lesen.
Ein Agent hat kein Konto.
Er handelt in jedem Repository, dessen Register ihn nennt, und nirgends sonst; auch ein Lauf in
der CI ist ein Agent mit einem Eintrag im Register.

Ob ein Repository öffentlich ist, gilt je Repository, und die Voreinstellung ist privat.
Ein öffentliches Repository wird ohne Anmeldung gelesen, ein privates von den Akteuren seines
Registers und von den Konten, denen der Namespace das Lesen gewährt.
Auf eine Anfrage ohne Leserecht antwortet ein privates Repository wie eines, das es nicht gibt.
Schreiben ist immer angemeldet.
Der Hub nimmt nur Repositories mit Register an: Der offene Modus bleibt ein lokaler Zustand, und
der erste Push muss ein Register bringen, das den Schlüssel der Anfrage als Owner führt.

\startinfobox
{\bf Eine Anmeldung ersetzt nie einen Beweis.}
Der Hub prüft dieselben Beweise mit der Operation des Kerns noch einmal und nimmt nichts an, nur
weil der Absender angemeldet ist.
Er darf mehr ablehnen als das Register, aber nie mehr erlauben.
Er schreibt keine Identität in ein Register und signiert nie für einen Akteur.
\type{op verify} kommt mit und ohne Hub zum selben Ergebnis, und ein übernommener Hub kann weder
eine Freigabe noch einen Eintrag des Operation Log fälschen.
\stopinfobox

\section[sec:ref-update]{Ref-Update und Aufzeichnung}

Lokal sagt das Operation Log, wer einen Zeiger bewegt hat.
Dieses Log verlässt den Klon nicht (ADR-011).
Für eine geteilte Referenz braucht es deshalb einen eigenen Beleg, den alle Leser prüfen können:
das signierte Ref-Update und die Aufzeichnung, in die der Hub es einträgt.

\subsection{Das Ref-Update}

Ein Ref-Update ist ein Beweis mit eigenem Typ, signiert mit dem Schlüssel aus dem Register.
Der signierte Text nennt die Adresse des Hubs, die ID, die der Hub dem Repository zugewiesen hat,
den öffentlichen Schlüssel und je Referenz:

\startitemize[packed]
\item den Namen, \type{heads/<Name>} oder \type{tags/<Name>};
\item \type{previous}, den Hash des letzten Eintrags der Aufzeichnung zu diesem Namen, oder
\type{-} für einen Namen, den der Hub nie gesehen hat;
\item \type{old} und \type{new}, je ein Commit oder \type{-}, sodass Anlegen und Löschen
dasselbe Ref-Update sind;
\item den Beleg, auf dem die Bewegung beruht: den Hash der Eröffnung eines Pull Requests, einer
Purge-Erklärung oder eines Owner-Override, sonst \type{-}.
\stopitemize

Der Text enthält keine Zeit und keinen Namen des Akteurs.
Der Akteur folgt aus Schlüssel und Register, und die Zeit setzt der Hub.
Ein Ref-Update kann mehrere Referenzen unter einer Signatur tragen; der Hub bewegt dann alle
oder keine.

\type{previous} tritt an die Stelle des erwarteten Commits im Compare-and-Swap.
Der Grund ist ein Fall, den der blosse Commit nicht unterscheidet: Eine Referenz geht von A nach
B und zurück nach A.
Ein altes, signiertes Ref-Update, das A erwartet, würde wieder passen und ein zweites Mal zählen.
Die Einträge zu einem Namen bilden aber eine Kette, und nach zwei Bewegungen ist der letzte
Eintrag ein anderer.
Ein Ref-Update zählt deshalb genau einmal, und die Kette eines Namens überlebt auch sein Löschen.

\subsection{Die Aufzeichnung}

Die Aufzeichnung der Ref-Updates ist eine Hash-Kette über alle angenommenen Ref-Updates eines
Repositorys, die der Hub allein schreibt.
Ein Eintrag hat zwei Teile und eine Signatur.

\startitemize[n]
\item Der signierte Text des Akteurs, Byte für Byte wie empfangen.
\item Die Zeilen des Hubs: eine laufende Nummer, die Zeit des Hubs, der Commit von \type{main},
dessen Register er verwendet hat, die Commits der Policies, je Referenz die Angabe, ob die
Bewegung ein Fast-Forward war, die Identität des Akteurs aus dem Register und der Hash des
vorherigen Eintrags.
\item Die Signatur des Hubs über den Hash des Eintrags.
\stopitemize

Der Hash eines Eintrags deckt beide Teile, die Signatur des Hubs steht ausserhalb.
Die Kette legt die Reihenfolge über alle Referenzen fest, wie die Hash-Kette des Operation Log
die Reihenfolge der Einträge festlegt (\in{Abschnitt}[sec:oplog-tip]).
Ein Eintrag entsteht genau dann, wenn sich Referenzen bewegen, und er wird vor der ersten
Bewegung geschrieben (Write-Ahead-Regel).
Abgelehnte Anfragen stehen nicht darin; ein Absender kann die Aufzeichnung also nicht füllen.

Die Antwort auf ein angenommenes Ref-Update ist der ganze Eintrag, die Quittung.
Der Klon prüft, dass sie seinen Text unverändert enthält und vom Hub signiert ist,
und bewegt erst dann seine Remote-Referenzen.
Jeder, der das Repository lesen darf, kann die Aufzeichnung über das Protokoll lesen und
nachrechnen.
Sie hält keinen Commit am Leben: Ein umgeschriebener Verlauf kann den Hub verlassen, und ein
Eintrag bleibt auch ohne seine Commits lesbar, weil der Hub bei der Annahme festhält, was sich
später nicht mehr berechnen lässt.

\subsection{Remote-Referenzen}

Im Klon ist eine Remote-Referenz eine Datei unter \type{refs/remotes/<Remote>/heads/} oder
\type{refs/remotes/<Remote>/tags/} mit zwei Zeilen: dem Commit und dem Eintrag der Aufzeichnung.
Nur \type{fetch} und \type{push} bewegen sie, mit je einem Eintrag im Operation Log.
Ein Fetch entfernt eine Remote-Referenz, die der Hub nicht mehr führt, im selben Eintrag;
\type{op restore} bringt sie zurück.
Die Tags eines Remote bleiben in dessen Namensraum, denn ein lokaler Tag entsteht nur durch
\type{tag} mit eigenem Eintrag.
Remote-Referenzen sind Wurzeln der Erreichbarkeit.

Gebaut ist davon, was ein Klon ablegt und liest; \type{fetch} und \type{push} folgen.
Die Konfiguration eines Remote ist eine Datei \type{.stemma/remotes/<Remote>} mit sechs Zeilen in
fester Reihenfolge: der Kopfzeile \type{stemma-remote 1}, der Adresse, der ID des Repositorys,
dem festgehaltenen Schlüssel des Hubs, einer zusätzlichen CA-Datei und der Angabe, ob der Hub
eine Anmeldung verlangt.
Der Name eines Remote ist eine einzelne Komponente nach den Regeln eines Branch-Namens.
Die zweite Zeile einer Remote-Referenz lautet \type{entry} und der Hash des Eintrags.
Für beide Dateien gilt die Regel der Objekte: Der Leser nimmt an, was der Schreiber schreibt,
und alles andere ist ein Schaden mit dem Pfad der Datei.
Mit den neuen Dateien steigt das Repository-Format auf Version 14 (ADR-009), und ein älteres
Programm lehnt ein solches Repository ab.

Wo ein Befehl eine Revision annimmt, nennt \type{remotes/<Remote>/heads/<Name>} oder
\type{remotes/<Remote>/tags/<Name>} eine Remote-Referenz.
Die Form sucht nie nach einem Branch oder Tag weiter, und das Präfix \type{remotes/} ist für neue
Branch- und Tag-Namen reserviert, wie \type{heads/} und \type{tags/}.
Die Prüfung der Referenzen prüft auch den Commit jeder Remote-Referenz.

\startinfobox
{\bf Purge und Remote-Referenzen.}
Der Plan eines Purge geht von den Remote-Referenzen aus wie von einem Branch.
Er bewegt und entfernt aber keine: Zeigt eine auf einen Commit, den der Purge umschreiben würde,
lehnt er ab und nennt jede solche Referenz.
Was ein Purge am Ende mit ihnen tut, ist offen (\in{Abschnitt}[sec:phase5-open]).
\stopinfobox

Ein veraltetes Ref-Update beantwortet der Hub mit dem Stand jeder betroffenen Referenz.
Der Klon meldet das und wiederholt nicht von selbst: Er müsste neu signieren, und zuerst
entscheidet der Akteur, ob er seinen Branch neu aufspielt oder einen Merge macht.

\section[sec:hub-enforcement]{Durchsetzung auf dem Hub}

Die lokale Durchsetzung hält einen Prozess nicht auf, der am Kern vorbei schreibt.
Der Hub schliesst diese Lücke für alles, was geteilt wird: Eine Referenz bewegt sich dort nur
mit einem signierten Ref-Update, und \type{main} bis auf zwei Ausnahmen nur, wenn das
Merge-Gate (\in{Abschnitt}[sec:merge-gate]) noch einmal gerechnet hat.
Die Ausnahmen sind der erste Push und der Purge; sie stehen bei den Garantien am Ende des
Abschnitts.

\subsection{Belege als signierte Erklärungen}

In Phase 4 liegen Pull Request, Freigaben und Prüfergebnisse im Operation Log.
Ein Hub sieht keinen Eintrag dieses Logs, denn es kann ein Geheimnis enthalten, bis ein Purge es
entfernt.
Die Belege bekommen deshalb eine eigene, tragbare Form: die signierte Erklärung.
Die Eröffnung eines Pull Requests, eine Freigabe, ein Prüfergebnis, das Schliessen, ein
Owner-Override und der Purge von \type{main} sind je ein Typ von Beweis.
Eine Erklärung bindet das Repository, den Pull Request und die Commits, aber keine Nummer, keine
Kette und keine Zeit eines Klons.

Der Hash der signierten Eröffnung ist der Name des Pull Requests.
Niemand vergibt ihn, er gilt in jedem Klon und auch ohne Hub, und ein übernommener Hub kann eine
Freigabe nicht einem anderen Pull Request zuordnen.
Die Erklärungen sind der Beleg, das Log hält die Handlung fest: Die Einträge \type{pr-open},
\type{pr-approve} und \type{check} nennen den Hash der Erklärung, die sie erzeugt haben.
Ein Klon und der Hub halten je eine Menge von Erklärungen und falten sie mit derselben Funktion.
Ein Fetch bringt die Erklärungen offener Pull Requests in den Klon, sodass das lokale Gate gleich
zählt wie der Hub.
Diese Form ersetzt die Einträge aus Phase 4 als Quelle des Gates und kostet eine neue Version des
Formats.

Der Hub prüft eine eingereichte Erklärung nach den Regeln des lokalen Schreibers: Signatur und
aktiver Schlüssel, eine Eröffnung nur durch den Schreibinhaber, eine Freigabe nur für den
aktuellen Kopf und nie die eigene.
Er fügt seine signierten Zeilen an, mit einer laufenden Nummer, der Zeit, dem Commit des
Registers und der Nummer des letzten Eintrags der Aufzeichnung,
und liefert die Erklärung allen Lesern des Repositorys aus.

\subsection{Wie sich main bewegt}

Der Klon baut den Merge-Commit, und der Hub rechnet ihn nach.

\startitemize[n]
\item Das Gate läuft im Klon wie bisher und baut den Merge-Commit $M$ auf dem Stand $T$ von
\type{main} auf dem Hub.
\item Der Klon lädt die Objekte hoch und sendet ein Ref-Update von $T$ nach $M$, das den Pull
Request nennt.
\item Der Hub prüft: $M$ hat den einzigen Elternteil $T$, und sein Tree ist der saubere Merge des
Kopfs auf $T$.
\item Er berechnet die Risikostufe aus der Policy von $T$, aus den Pfaden der Änderung und aus
der Eskalation bei abgeschwächten Tests (\in{Abschnitt}[sec:policy-file]).
\item Er faltet die Erklärungen und zählt Freigaben und Prüfergebnisse gegen das Register von
$T$.
\item Er prüft die Schranken des Gates: die Pause bei roter Spitze, die Reservierung und das
Fenster ungeprüfter Commits.
\item Er schreibt den Eintrag der Aufzeichnung, bewegt \type{main} und antwortet mit der Quittung.
\item Erst nach der Quittung bewegt der Klon sein eigenes \type{main}.
\stopitemize

Das Ref-Update nennt keine Stufe, und nichts, was der Absender behauptet, geht in das Urteil ein.
Ein $M$, das die Policy ändert, wird nach der alten Policy beurteilt.
Der Hub führt keinen Merge aus und erzeugt keinen Commit (ADR-015); er führt auch keinen fremden
Code aus.
Ein Prüfergebnis ist die Erklärung eines Akteurs aus dem Register, und niemand beweist den Lauf.

Eine Ablehnung hinterlässt nur unbenannte Objekte und nichts, was zurückzunehmen wäre.
Hat sich \type{main} inzwischen bewegt, holt \type{pr merge} den neuen Stand, rechnet neu und
signiert neu, eine begrenzte Zahl von Malen.
Ein gewöhnlicher Push wiederholt dagegen nicht von selbst
(\in{Abschnitt}[sec:ref-update]).
Ohne Verbindung bleiben Commits, Branches, das Eröffnen eines Pull Requests, Freigaben und
Prüfungen möglich; nur die Bewegung des geteilten \type{main} wartet auf den Hub.

Der Zustand \quotation{offen, obwohl \type{main} die Änderung enthält}, den das lokale Gate
reparieren muss (\in{Abschnitt}[sec:merge-gate]), entsteht auf dem Hub nicht.
Der Zustand \type{merged} ist dort keine Erklärung: Er folgt aus dem Eintrag der Aufzeichnung,
der den Pull Request nennt, und Bewegung und Merge sind ein Schritt unter einer Sperre.

\subsection{Branches, Tags und Purge}

Der Schreibinhaber eines ungeschützten Branch ist auf dem Hub der Akteur des letzten Eintrags der
Aufzeichnung zu diesem Namen.
Die Lease läuft ab der Zeit des Hubs in diesem Eintrag, mit der Dauer aus der Policy.
Der Hub braucht dafür keinen eigenen Zustand, und jeder Leser kann es nachrechnen.
Der Schreibinhaber darf seinen Branch beliebig bewegen; die Aufzeichnung hält fest, wenn eine
Bewegung kein Fast-Forward war.
Ein Owner-Override ist eine signierte Erklärung mit Begründung, die alle Leser sehen.

Ein Tag wird nach derselben Regel gesetzt und gelöscht wie lokal.
\type{main} bewegt sich nach dem ersten Push nur von $T$ nach $M$, mit einer Ausnahme:
Für einen Purge legt ein menschlicher Owner nach dem Register von $T$ eine signierte
Purge-Erklärung vor, mit alter Spitze, neuer Spitze und Begründung.
Der Eintrag ist für alle Leser sichtbar, und Erklärungen über entfernte Commits zählen nicht
mehr.

\subsection{Garantien und Grenzen}

Auf dem Hub gilt: Jede Bewegung einer Referenz ist von einem Schlüssel signiert, der im Register
an der Spitze von \type{main} aktiv war, und \type{main} bewegt sich von $T$ nach $M$ nur, wenn
das Gate bestanden ist, das der Hub selbst gerechnet hat.
Ohne Gate kommen zwei Stände nach \type{main}: der erste Push in ein leeres Repository, dessen
Register den Schlüssel der Anfrage als Owner führen muss, und der Purge eines menschlichen
Owners.
Eine Ablehnung nennt dem Absender den vollen Grund: die Rolle, den Inhaber der Lease, die
berechnete Stufe, die gezählten und fehlenden Freigaben und Prüfergebnisse oder die Schranke.
Sie verrät nichts, denn wer sie bekommt, darf lesen und könnte alles selbst berechnen.
Ablehnungen stehen in keinem geteilten Format.
Sichtbar ist, worauf das Urteil beruht: die Erklärungen, die Aufzeichnung und das Repository.

\startinfobox
{\bf Eine Operation, zwei Eingänge.}
Der Kern nimmt ein Ref-Update oder eine Erklärung an und nimmt die Sperre der Referenzen dabei
selbst, wie das Gate.
Die beiden Eingänge heissen \type{accept_ref_update} und \type{accept_statement}.
Vor der Sperre laufen nur reine Berechnungen über unveränderliche Objekte: Form und Signatur, der
Tree-Merge, die Stufe.
Unter der Sperre steht, was von veränderlichem Zustand abhängt: \type{previous}, der Zusammenhang
der Objekte, Register, Erklärungen, Schranken, der Eintrag und die Bewegung.
Der Upload läuft ohne Sperre.
\stopinfobox

\section[sec:hub-anchor]{Anker und der Schlüssel des Hubs}

Signaturen verhindern nicht, dass jemand das Ende seines Operation Log abschneidet.
Das findet nur ein Anker: Nummer und Hash eines Eintrags, aufbewahrt von jemandem, der die Kette
nicht selbst schreibt (\in{Abschnitt}[sec:oplog-tip]).
Bisher war das eine Kopie, die ein Mensch aufbewahrt.
Mit dem Hub gibt es den Anker in beide Richtungen.

Der Hub hat einen eigenen Ed25519-Schlüssel je Instanz.
Er signiert damit nur seine eigenen Aussagen: jeden Eintrag der Aufzeichnung, seine Zeilen an
einer Erklärung und die Bestätigung eines Ankers.
Für einen Akteur spricht dieser Schlüssel nie.
Ein Klon hält den Schlüssel beim ersten Kontakt in der Konfiguration des Remote fest und zeigt den
Fingerabdruck; ein anderer Schlüssel ist später ein Fehler, den ein Mensch auflöst.
Quittung und Bestätigung sind damit Beweise, die ein Klon prüft, ohne dem Transport zu vertrauen.
Zwei signierte Einträge mit derselben Nummer und verschiedenen Hashes beweisen einen Hub, der
seine Aufzeichnung umschreibt.

\startitemize
\item Die Quittung ist ein Anker der Aufzeichnung, den der Klon hält und der Hub nicht ändern
kann.
\item Vom Operation Log bestätigt der Hub allein den Anker der Spitze.
Der Klon sendet Nummer und Hash seines letzten Eintrags, signiert vom Akteur; der Hub bewahrt
den Anker auf und antwortet mit einer signierten Bestätigung, die seine Zeit und die Nummer des
letzten Eintrags der Aufzeichnung nennt.
\stopitemize

Der Hub sieht dabei keinen Eintrag und kein Geheimnis.
Die Aussage aus ADR-011 lautet genauer: Ein Hub sieht nie einen Eintrag des Logs, er bestätigt
Nummer und Hash der Spitze.
Ein Push sendet den Anker nach der Quittung, als eigene Anfrage.
Der Anker deckt so den Eintrag \type{push}, der die Quittung nennt, und Log und Aufzeichnung
binden einander.
Ein Log gehört zu einem \type{.stemma} und trägt eine zufällige Log-ID, denn ein Akteur kann zwei
Klone mit demselben Schlüssel haben, und mehrere Akteure schreiben in ein Log.

Einen widersprechenden Anker, dieselbe Nummer mit anderem Hash oder eine kleinere Nummer, hält
der Hub als solchen fest, mit der Signatur des Akteurs als Beleg.
Er bestätigt ihn nicht und blockiert nichts.
Alle Leser des Repositorys sehen ihn; was daraus folgt, entscheiden Menschen.
Ein wiederhergestelltes Backup hält so keinen Akteur auf, dessen Beweise gültig sind.

\subsection{Die Ordnung über Klone hinweg}

Das Register schreibt die Grenze eines ausgemusterten oder widerrufenen Schlüssels als Nummer des
lokalen Logs (\in{Abschnitt}[sec:register]).
Jeder Klon hat eigene Nummern, die Grenze gilt also nur in dem Log, das sie geschrieben hat.
Über Klone hinweg gilt die Ordnung des Hubs.
Ein Widerruf wirkt ab dem Eintrag der Aufzeichnung, der \type{main} auf den Commit mit dem
Widerruf bewegt.
Eine Erklärung liegt davor, wenn ihre Zeilen des Hubs eine kleinere Nummer nennen.
Ein Eintrag im Log eines anderen Klons liegt nur dann davor, wenn ein Anker ihn deckt, dessen
Bestätigung eine kleinere Nummer nennt.
Der Hub entscheidet das mit Nummern unter einer Sperre und nie mit seiner Uhr.

Eine Kopie der Policy hält der Hub nicht.
Er liest sie aus dem Basis-Commit, wie der Kern es lokal tut; es gibt also nichts, was er
signieren müsste.

\section[sec:namespaces]{Namespaces}

Ein Kurs braucht Regeln, die für alle Repositories der Studierenden gelten, und eine Lehrperson,
die in jedem davon eingreifen darf.
Dafür gibt es den Namespace: eine flache Gruppe von Repositories mit einer ID, einem Namen und
genau einer Namespace-Policy.
Eine Organisation, ein Kurs und ein Projekt sind dasselbe und kein eigener Typ.
Die Adresse eines Repositorys hat immer zwei Teile, Namespace und Repository.

Die Regeln eines Namespace liegen in einem eigenen Stemma-Repository, dem Namespace-Repository.
Es enthält die Namespace-Policy als Datei und ein Register; dessen menschliche Owner sind die
Owner des Namespace.
Eine Änderung nimmt den gewöhnlichen Weg: ein Pull Request der Stufe 3, mit Erklärungen und einem
Eintrag der Aufzeichnung.
Ein Konto reserviert den Namen, und der Namespace entsteht mit dem ersten Push seines
Namespace-Repositorys.

Für ein Repository werden beide Policies ausgewertet, und die strengere zählt.
Die annehmende Operation liest die Policy von $T$ und die Namespace-Policy an der Spitze von
\type{main} des Namespace-Repositorys und nimmt je Änderung das Maximum: die Stufe über die
berührten Pfade, die Zahl der Freigaben, die engeren Schranken.
Eine Namespace-Policy kann deshalb nur anheben, und die Mindeststufen des Kerns bleiben, weil ein
Maximum nie unter sie fällt.
Die Datei der Namespace-Policy ist selbst immer Stufe 3.
Ein Fetch bringt sie in den Klon, damit das lokale Gate zum selben Urteil kommt.

Das Register bleibt die einzige Quelle der Rollen in einem Repository.
Owner eines Namespace haben Rechte nur oberhalb des Repositorys: anlegen, umbenennen, löschen,
Leser bestimmen.
Die Namespace-Policy kann aber mit einer Zeile verlangen, dass das Register jedes Repositorys die
Owner des Namespace als Owner führt.
Der Hub lehnt dann einen ersten Push und jede Bewegung von \type{main} ab, deren Register sie
nicht enthält.
So kann ein Kurs den Owner-Override der Lehrperson garantieren, und \type{op verify} kommt ohne
Hub zum selben Ergebnis, weil die Rolle im Register steht.

\startitemize
\item Ein Owner des Namespace gewährt einem Konto \type{create} (Repositories anlegen) oder
\type{read} (alle privaten Repositories des Namespace lesen), mit einer signierten Erklärung.
\item Wer ein Repository lesen darf, liest alle seine Branches und Tags, die Aufzeichnung, die
Erklärungen und die Anker.
\item Einen Owner-Override sehen die Leser des Repositorys.
Dazu liefert der Hub je Namespace eine Liste aller Owner-Overrides, gefiltert auf das, was der
Anfragende lesen darf: Die Lehrperson sieht den ganzen Kurs, Studierende sehen ihre eigenen
Repositories.
\item Ein Owner des Namespace kann ein Repository einfrieren.
Der Hub lehnt dann jedes Ref-Update und jede Erklärung ab; das Lesen geht weiter.
Zurücksetzen kann es keine Rolle des Namespace: Der Hub würde sonst eine Bewegung erlauben, die
das Register nicht deckt.
Der Weg zurück sind gewöhnliche Beweise, ein Widerruf und ein Revert oder ein Purge.
\item Löschen geschieht in zwei Schritten: Nach der signierten Erklärung eines Owners des
Namespace liefert der Hub das Repository nicht mehr aus, und erst nach einer Frist, in der ein
Owner es wiederherstellen kann, wird es entfernt.
\stopitemize

\startinfobox
{\bf Namen und IDs.}
Ein Name besteht aus \type{a-z}, \type{0-9} und dem Bindestrich im Innern, hat 1 bis 64 Zeichen
und beginnt mit einem Buchstaben oder einer Ziffer.
Die ID eines Repositorys sind 128 zufällige Bits, die der Hub beim Anlegen vergibt und nie
wiederverwendet.
Ein Ref-Update bindet die ID und nicht den Namen: Nach dem Umbenennen bleiben die Aufzeichnung
und vorbereitete Ref-Updates gültig, und der alte Name antwortet mit dem neuen.
\stopinfobox

\section[sec:threat-model]{Was ein Hub nicht verhindert}

Agenten laufen mit den Rechten ihres Menschen auf dessen Maschine.
Ein bösartiger Prozess mit denselben Rechten des Betriebssystems erreicht alles, was dieser
Benutzer erreicht: die Dateien des Repositorys samt \type{.stemma}, jede Schlüsseldatei und jeden
\type{ssh-agent} des Benutzers, laufende Agenten, das Programm \type{stemma} und seine Anzeige.
Nicht erreichbar sind der Hub und sein Schlüssel, Schlüssel auf anderen Maschinen und eine
Signatur, die eine Berührung eines Hardware-Tokens braucht.
\in{Tabelle}[tab:threats] ordnet ein, was der Hub gegen einen solchen Prozess ausrichtet.

\startplacetable[title={Was der Hub gegen einen Prozess mit den Rechten des Akteurs ausrichtet},
reference=tab:threats]
\startxtable[frame=off,option=stretch]
\startxrow[topframe=on,bottomframe=on]
\startxcell[width=0.2\textwidth] {\bf Wirkung} \stopxcell
\startxcell {\bf Angriff} \stopxcell
\stopxrow
\startxrow
\startxcell Verhindert \stopxcell
\startxcell Eine Referenz, ein Objekt oder ein Eintrag, am Kern vorbei und ohne Schlüssel
geschrieben, erreicht den Hub nie.
Eine gefälschte Freigabe ohne Schlüssel ist kein Beweis.
Die Aufzeichnung lässt sich nicht fälschen, weil der Prozess den Schlüssel des Hubs nicht
erreicht und die Quittungen in den Klonen liegen. \stopxcell
\stopxrow
\startxrow
\startxcell Sichtbar gemacht \stopxcell
\startxcell Ein abgeschnittenes oder umgeschriebenes Operation Log zeigt sich als
widersprechender Anker. \stopxcell
\stopxrow
\startxrow
\startxcell Ohne Gewinn \stopxcell
\startxcell Das Token sagt nur, welcher Schlüssel spricht. \stopxcell
\stopxrow
\startxrow[bottomframe=on]
\startxcell Offen \stopxcell
\startxcell Der Prozess benutzt den Schlüssel des Akteurs.
Für den Hub ist er dann der Akteur. \stopxcell
\stopxrow
\stopxtable
\stopplacetable

Gegen den benutzten Schlüssel baut Phase 5 einen zweiten Schlüsseltyp und einen Bericht.
Das Register nimmt neben Ed25519 den Hardware-Schlüssel an
(\type{sk-ssh-ed25519@openssh.com}; mit einem echten Hardware-Token ist der Weg noch nicht
erprobt, ADR-020), und der Kern verlangt bei jeder Signatur eines solchen
Schlüssels das Merkmal \quotation{Benutzer anwesend}, ohne Ausnahme nach Typ des Beweises.
Der Schlüsseltyp ist freiwillig und im Register sichtbar: Jeder Leser sieht, welcher Owner wie
geschützt ist.
Der Preis ist eine Berührung je Signatur.
Ein Push eines Menschen mit Hardware-Schlüssel kostet drei, für die Anmeldung, das Ref-Update und
den Anker; ein Fetch eines privaten Repositorys kostet eine.
Wer oft pusht, tut es als Agent mit einem Schlüssel in einer Datei, und der Mensch mit dem
Hardware-Schlüssel gibt frei und führt die Handlungen der Stufe 3 aus.

Der Bericht ist ein Befehl, der aus der Aufzeichnung und den Erklärungen zusammenstellt, was ein
Schlüssel in einem Bereich von Nummern getan hat: welche Referenzen er bewegt hat, welche
Erklärungen von ihm stammen und welche Merges nach \type{main} darauf beruhten.
Nach einem Diebstahl wirkt der Widerruf ab dem Eintrag, der \type{main} auf ihn bewegt, und die
Freigaben des Schlüssels verfallen.
Das Aufräumen bleibt bei den Menschen, mit einem Revert oder einem Purge.

\startwarnbox
{\bf Grenzen, die bleiben.}
\startitemize[packed]
\item Ein Schlüssel als Datei oder in einem \type{ssh-agent} ist der Akteur für jeden Prozess des
Benutzers.
Eine Passphrase schützt nur die ruhende Datei, und kein Prüfer sieht sie.
\item Die Schlüssel mehrerer Akteure unter einem Benutzer fallen zusammen.
Ein Prozess mit den Schlüsseln zweier Agenten eröffnet mit dem einen und gibt mit dem anderen
frei; er bringt eine Änderung der Stufe 0 oder 1 allein nach \type{main}.
Sichtbar bleibt es in der Aufzeichnung und den Erklärungen.
\item Das Programm und seine Anzeige liegen in Reichweite: Einem Menschen kann das eine gezeigt
und das andere zur Signatur vorgelegt werden.
\item Ein Klon kann seine Konfiguration nicht vor seinem Benutzer schützen.
Nach einer Kompromittierung führt der Weg zurück über einen frischen Klon,
der mit der Aufzeichnung verglichen wird.
\item Mit dem gestohlenen Schlüssel eines Owners ohne Hardware-Token übernimmt ein Prozess das
Register und damit die Regeln des Repositorys auf dem Hub.
Die Übernahme bleibt sichtbar, und kein Klon verliert seine Objekte.
Der Weg zurück liegt oberhalb des Repositorys: Der Betreiber sperrt den Schlüssel, und ein Owner
des Namespace friert das Repository ein.
\stopitemize
\stopwarnbox

Daraus folgen Empfehlungen: ein Hardware-Schlüssel für jeden Owner, zwei Owner-Schlüssel an
getrennten Orten, und die Schlüssel von Agenten getrennt von denen des Menschen, unter einem
eigenen Benutzer des Betriebssystems oder als Geheimnis der CI.
Eine Passphrase und \type{ssh-add -c} sind sinnvoll, beweisen aber niemandem etwas.

\section[sec:no-peer]{Zwei Klone ohne Hub}

Register und Signaturen liegen im Repository, deshalb prüft jeder Klon einen einzelnen Beweis
ohne Hub.
Die Durchsetzung auf einem geteilten Branch braucht zusätzlich eine Ordnung, die jemand schreibt,
der nicht der Akteur ist: für das Ausmustern eines Schlüssels, die Schranken des Gates, die
Schreibhoheit und die Namespace-Policy im Augenblick eines Merges.
Zwei Klone haben je ein Log mit eigenen Nummern und je ein eigenes \type{main}.
Ein Klon, der die Ordnung für andere festlegt, wäre ein Hub unter anderem Namen.

Ein Klon ist deshalb nie die Autorität über einen geteilten Branch und nimmt nie einen Push an.
Zwei Klone lesen voneinander; geschrieben wird auf einen Hub.
Der Transport von Objekten und Erklärungen zwischen Klonen bleibt möglich, weil Objekte über
ihren Hash adressiert sind und Erklärungen ohne Hub gültig sind.
Gebaut wird er nicht; der Auslöser wäre ein Austausch ohne erreichbaren Hub, etwa in einem
Schulnetz ohne Internet.
Auch dann wäre der andere Klon nur ein Bote: Nichts von ihm zählt, bevor der Hub es hat.

Im offenen Modus arbeitet man allein; das Teilen beginnt mit dem Register.
Wer ohne den zentralen Hub arbeiten will, braucht einen eigenen Hub.
Die Autorität ist eine Rolle und kein bestimmter Server: Ein Klon kennt einen Hub an seiner
Adresse und seinem Schlüssel, und ein neues Repository auf einem Hub entsteht jederzeit aus einem
sauberen Klon.
Geht ein Hub verloren, geht die Aufzeichnung verloren und kein Objekt.

\section[sec:phase5-open]{Offene Fragen}

\startitemize
\item Die Befehle: Namen, Optionen und Exit-Codes von \type{remote}, \type{clone}, \type{fetch}
und \type{push}, der Befehle für Erklärungen und Anker und der Befehle eines Namespace.
Jeder Teil entsteht mit seiner Umsetzung (\in{Kapitel}[chap:cli]).
\item Purge auf dem Hub: Die Regel der Annahme steht.
Offen ist, wie Klone einem \type{main} folgen, das ein Purge auf dem Hub umgeschrieben hat,
wie ein Purge Remote-Referenzen behandelt (bis dahin lehnt er ab, wenn eine auf einen Commit
zeigt, den er umschreiben würde)
und was aus den Erklärungen eines Klons über entfernte Commits wird (ADR-014).
\item Messung: Die Schwellen für ein Pack-Format, die Haltezeit der Sperre der Referenzen auf dem
Hub und die Zahl der Wiederholungen eines Merges wegen \type{Stale} sind nicht gemessen
(ADR-015, ADR-023).
\item Einträge aus der Zeit vor der Hash-Kette des Operation Log deckt kein Anker des Hubs.
\item Ein Repository kann seinen Namespace nicht wechseln.
\item Wie eine Änderung der Stufe 3 in einem Namespace-Repository mit einem einzigen menschlichen
Owner freigegeben wird, etwa in einem Kurs mit einer Lehrperson.
\item Strengere Mindeststufen für Firmenkunden bleiben ausserhalb des Umfangs; ADR-027 hält fest,
was das Modell dafür offen lässt.
\item Ob und wie der Hub an Dritte weitergegeben wird.
Davon hängt ab, was \quotation{ein eigener Hub} für sie bedeutet.
\stopitemize

Commits und Merges auf dem Hub für Agenten ohne Klon baut Phase 5 nicht; die Richtung dafür steht
in \in{Kapitel}[chap:phase6].
